% Claim proof file for row claim_3_1 -- "CDMGRestrictions".
% Auto-generated stub by scaffold/solve_chapter.py at row start. A manager
% agent dedicated to writing the TeX proof fills in the proof body below.
%
% This file is a sibling of `main.tex` inside `<subsection>/tex/`. The
% `subfiles` package lets it render both standalone and as part of
% `main.tex` via `\subfile{...}`.
%
% The statement is restated above the proof so the file is self-contained
% when rendered on its own. The proof must:
%   - Stay close to the lecture notes' proof if one exists (always search
%     the LN first for a `\begin{proof}` block immediately following the
%     claim; copy / lightly edit when present).
%   - Otherwise use the LN paradigm and cite prior claims by their `ref`.
%   - Be self-contained enough that a mathematician can follow it without
%     re-reading the LN.
%   - Have no `sorry`, `omitted`, or "obvious" shortcuts.

\documentclass[main]{subfiles}

\begin{document}
% ref: claim_3_1 (proof, with statement restated for readability)
% title: CDMGRestrictions
\def\rowref{claim\_3\_1}
\phantomsection\label{claim_3_1}

% --- statement (restated) ---------------------------------------------------
\begin{Rem}[CDMGRestrictions]
    Let $G = (J, V, E, L)$ be a CDMG (def \ref{def-cdmg}); we use the notation of \ref{not-cdmg} throughout. Recall, per def \ref{def-cdmg}, that $E \subseteq (J \cup V) \times V$, that $L \subseteq V \times V$ with $L$ symmetric and irreflexive, and that $J \cap V = \emptyset$.

    The restrictions encoded in def \ref{def-cdmg} have the following three consequences. In the LN wording, the clauses ``$j \hus v \notin G$'' and ``edges $j \tuh v$ are allowed'' carry an implicit universal quantification over the second-position variable $v$; we surface these quantifiers explicitly below, with $v$ ranging over $J \cup V$ in the first clause and over $V$ only in the second (see also the addition spelled out under sub-claim ii.).

    \begin{enumerate}[label=\roman*.)]
        \item \emph{No edge with an arrowhead at a $J$-node.}
              For every $j \in J$ and every $v \in J \cup V$,
              \[ j \hus v \notin G. \]
              Unfolding the shorthand $\hus$ via def \ref{not-cdmg} item~6 (so $j \hus v \in G$ abbreviates $(v, j) \in E \lor (j, v) \in L$), this is equivalent to the set-theoretic conjunction
              \[ (v, j) \notin E \;\land\; (j, v) \notin L. \]

        \item \emph{The type of $G$ does not forbid a directed edge from a $J$-node into a $V$-node.}
              For every $j \in J$ and every $v \in V$, the ordered pair $(j, v)$ satisfies $j \in J \cup V$ and $v \in V$, so the typing constraint $E \subseteq (J \cup V) \times V$ of def \ref{def-cdmg} is compatible with $(j, v) \in E$; equivalently, no restriction in def \ref{def-cdmg} forbids the inclusion of $(j, v)$ in $E$ (in the notation of \ref{not-cdmg}: no restriction forbids $j \tuh v \in G$). This is a \emph{permission} statement, not an existence statement: $(j, v) \in E$ is admissible by the definition for every $(j, v) \in J \times V$, but is not asserted to hold.

        \item \emph{No two $J$-nodes are adjacent.}
              For every $j_1 \in J$ and every $j_2 \in J$, the nodes $j_1$ and $j_2$ are \emph{not} adjacent in $G$ in the sense of def \ref{def-edge-relations} item~i.; equivalently,
              \[ (j_1, j_2) \notin E \;\land\; (j_2, j_1) \notin E \;\land\; (j_1, j_2) \notin L. \]
    \end{enumerate}
\end{Rem}

% --- proof ------------------------------------------------------------------
\begin{proof}
All three sub-claims rest on the same underlying tool: the typing constraints carried by def \ref{def-cdmg}, namely $E \subseteq (J \cup V) \times V$, $L \subseteq V \times V$, and $J \cap V = \emptyset$. We prove each in turn.

\medskip
\emph{Proof of (i).}
Fix $j \in J$ and $v \in J \cup V$, and suppose, for contradiction, that $j \hus v \in G$. By def \ref{not-cdmg} item~6, the shorthand $j \hus v \in G$ abbreviates the disjunction
\[ (v, j) \in E \;\lor\; (j, v) \in L. \]
We rule out both disjuncts.
\begin{itemize}
    \item If $(v, j) \in E$, then by the typing constraint $E \subseteq (J \cup V) \times V$ of def \ref{def-cdmg}, the second component of any element of $E$ lies in $V$; hence $j \in V$. But $j \in J$ by assumption and $J \cap V = \emptyset$ (def \ref{def-cdmg}), so $j \in J \cap V = \emptyset$, a contradiction.
    \item If $(j, v) \in L$, then by the typing constraint $L \subseteq V \times V$ of def \ref{def-cdmg}, the first component of any element of $L$ lies in $V$; hence $j \in V$. The same contradiction $j \in J \cap V = \emptyset$ follows.
\end{itemize}
Both disjuncts therefore fail, so $j \hus v \notin G$. Unfolding the shorthand once more, this is the conjunction $(v, j) \notin E \land (j, v) \notin L$, as stated.

\medskip
\emph{Proof of (ii).}
Fix $j \in J$ and $v \in V$. We must show that the typing constraint $E \subseteq (J \cup V) \times V$ of def \ref{def-cdmg} is compatible with the ordered pair $(j, v)$ -- i.e.\ that $(j, v) \in (J \cup V) \times V$, so that nothing in def \ref{def-cdmg} forbids the membership $(j, v) \in E$. The membership $(j, v) \in (J \cup V) \times V$ is, by the definition of a Cartesian product, the conjunction
\[ j \in J \cup V \;\land\; v \in V. \]
The second conjunct is the standing hypothesis $v \in V$. For the first conjunct, $j \in J$ gives $j \in J \cup V$ by inclusion of the left summand into the union. Hence $(j, v) \in (J \cup V) \times V$, which is the typing precondition for an element of $E$. No further restriction in def \ref{def-cdmg} excludes the pair $(j, v)$ from $E$ (the only other constraints concern $L$, not $E$, and the disjointness $J \cap V = \emptyset$ is compatible with $(j, v)$ since $j \in J$ and $v \in V$ are unrelated by intersection). Thus the inclusion $(j, v) \in E$ is admissible by the definition, which is the permission statement claimed. (We emphasise that this argument does not assert that $(j, v) \in E$ holds for any particular CDMG: it asserts only that the definition of CDMG does not preclude it.)

\medskip
\emph{Proof of (iii).}
Fix $j_1 \in J$ and $j_2 \in J$, and suppose, for contradiction, that $j_1$ and $j_2$ are adjacent in $G$. By def \ref{def-edge-relations} item~i., adjacency unfolds to the disjunction
\[ (j_1, j_2) \in E \;\lor\; (j_2, j_1) \in E \;\lor\; (j_1, j_2) \in L. \]
We rule out all three disjuncts.
\begin{itemize}
    \item If $(j_1, j_2) \in E$, then $E \subseteq (J \cup V) \times V$ (def \ref{def-cdmg}) forces the second component $j_2 \in V$. But $j_2 \in J$, so $j_2 \in J \cap V = \emptyset$ (def \ref{def-cdmg}), a contradiction.
    \item If $(j_2, j_1) \in E$, the same constraint applied to the second component yields $j_1 \in V$; combined with $j_1 \in J$ this gives $j_1 \in J \cap V = \emptyset$, a contradiction.
    \item If $(j_1, j_2) \in L$, then $L \subseteq V \times V$ (def \ref{def-cdmg}) forces both components into $V$; in particular $j_1 \in V$, hence $j_1 \in J \cap V = \emptyset$, a contradiction. (The same step also yields $j_2 \in V$; either component on its own suffices.)
\end{itemize}
All three disjuncts therefore fail, so $j_1$ and $j_2$ are not adjacent; equivalently the conjunction $(j_1, j_2) \notin E \land (j_2, j_1) \notin E \land (j_1, j_2) \notin L$ holds, as stated.

The argument above does not separate the cases $j_1 = j_2$ and $j_1 \neq j_2$: the typing-driven contradictions apply uniformly. We note in passing that the diagonal case $j_1 = j_2 =: j$ also follows: the $E$-branches give $j \in V$ as above, and the $L$-branch additionally contradicts the irreflexivity of $L$ from def \ref{def-cdmg} (which forbids $(j, j) \in L$) on top of forcing $j \in V$. No diagonal subtlety arises.
\end{proof}

\end{document}
